\chapter{Interference, allocation policies and software adoption}
\label{v13:sw-networks}

The preceding relationship analysis showed why a policy may require joint
graph and action information. This chapter examines the allocation problem
itself. It distinguishes network and grouped targets, develops the two initial
software references, and compares spatial, experimental, learned-policy and
simulation alternatives. The recurring small graphs and cost denominators are
matched comparison devices; each implementation's different assumptions and
failure conditions remain part of the argument.

\section{The unit receiving treatment need not be the independent unit}

A household member's offer can influence another member's borrowing, and an
outreach campaign can affect connected firms. In such settings a potential
outcome depends on more than the individual's own action. Writing only
$Y_i(a_i)$ hides the policy being evaluated. A more explicit representation is
$Y_i(a_i,\mathbf a_{-i})$, often restricted through an exposure mapping of
other units' actions. The mapping and the network are scientific inputs;
neither is established merely by finding matching customer identifiers.

Two broad software constructions appear in the inspected sources. MossSpider
uses a network and exposure summaries in a network-targeted estimation
workflow. Inferference uses grouped formulas, allocation probabilities and
inverse-probability weighting. These constructions have different observation
laws and potentially different targets. Their package names do not establish
interchangeability. \cite{v13_sw_mossspider,v13_sw_inferference}

An allocation policy illustrates the target. Suppose two linked borrowers have
binary outreach actions $A_1,A_2$, independently assigned with probability $p$
under the proposed policy. Stipulate the synthetic financial outcome
\[
 Y_i(a_i,a_j)=100a_i+40a_j-30a_i a_j
\]
in CNY. The first term represents the own-action contribution, the second a
spillover and the last diminishing combined benefit. These are hypothetical
structural values, not bank estimates. Averaging under the allocation gives
\[
 V(p)=140p-30p^2.
\]
For $p=0.2$ the mean is 26.8 CNY; for $p=0.8$ it is 92.8 CNY. The allocation
contrast is 66 CNY per borrower. Exact enumeration of the four assignments
checks the formula. \cite{v13_project_exact}

The direct effect of changing only the individual's own action is instead
$100-30a_j$, or $100-30p$ when the neighbor is independently treated with
probability $p$. This differs from the total allocation contrast. Reporting
one as the other can misprice the campaign even when every arithmetic
calculation is correct. The bank must specify whether it changes one action,
the neighborhood's allocation or an entire branch policy, and which costs
are included in the outcome.

\section{MossSpider: network semantics are part of the estimand}

The inspected \texttt{NetworkTMLE} documentation describes expected incidence
under a treatment plan, with Monte Carlo integration over allocations for
stochastic plans. It states a direct-contact interference assumption. The
constructor rejects self-loops, noninteger node identities and nonbinary
exposure, creates a relabeled network while retaining original identities,
and constructs an unweighted adjacency matrix. \cite{v13_sw_mossspider}

These details matter for bank graphs. A weighted guarantee network is not the
same as an unweighted adjacency graph. Removing weights can turn a small
guarantee into the same exposure contribution as a large one. Relabeling nodes
is harmless only when results can be mapped back to the correct identities.
A linkage error can attach the correct numerical estimate to the wrong firm.
Semantic conformance therefore needs an identity roundtrip and an explicit
policy for unsupported weighted relationships.

Direction requires particular review. The documentation initially describes an
undirected graph but also explains a directed influence convention opposite
the graph arrow, and the selected constructor includes directed handling. This
combination requires a concrete test with an asymmetric small graph. One should
not infer direction from a diagram's appearance or assume that every library
uses source-to-target influence. A two-node one-way example with distinguishable
covariates can determine which neighbor contributes to each exposure summary.
Section~\ref{sec:exec24-interference-reuse} develops the selected mapping,
policy and variance paths. Their inspected direction convention still needs
matched native tests; directed bank-network adoption remains a proposed use. \cite{v13_sw_mossspider}

Degree restrictions also change the target. The documentation treats nodes
outside a requested degree range as background, with their exposures unchanged
by the intervention. This can be scientifically useful, but it is not the
same as deleting those nodes from the graph or intervening on all customers.
The response must state which nodes are in the target population, which remain
background and how their factual exposure affects the intervention outcome.

The source disclaims a missing-data procedure. Missing node attributes,
uncertain links and omitted customers are common in banking, but the existence
of an exposure-mapping function does not solve those observation problems.
One must either supply defensible complete data in the declared regime or
develop a separate missingness/linkage construction. Network-targeted estimation
does not automatically qualify ownership or inclusion truth.

The observed MIT text makes a restricted adapter worth considering. Full adoption
requires qualification of the exposure-model and mapping code, targeting step,
Monte Carlo integration and variance calculation, together with actual
dependency/runtime checks. The developed study retains deterministic-policy
refusals, background-node treatment and a potentially negative variance
quadratic; these restrictions belong in any proposed adapter. A claim of general contagion, dynamic financial-network equilibrium or
arbitrary spillovers would go beyond the selected implementation evidence.

\section{Inferference: grouped allocation and numerical integration}

The inspected \texttt{interference} function uses a formula that identifies
outcome, exposure, optional participation, propensity covariates and a grouping
variable. It accepts allocation probabilities and propensity-model choices
including \texttt{glm}, \texttt{glmer}, \texttt{HLfit} and an oracle route.
Its documented causal-estimation path is IPW with naive or robust variance
options. The selected code sorts records by group and checks group count,
allocation range and the variance option. \cite{v13_sw_inferference}

A group-based model can be appropriate when the scientific design justifies
within-group interference and the required between-group sampling assumptions.
It does not become a model of arbitrary connected borrowers by assigning
everyone a group label. Cross-group guarantees or shared market shocks can
violate the intended observation law. The grouping decision needs domain
evidence and should be subjected to sensitivity analysis where boundaries
are uncertain.

The default logit integrand is restricted to a single random effect in the
documented implementation. The code rejects more than one when that integrand
is selected. Thus a proposal with household and branch random effects needs a
different justified integration path; silently dropping one random effect is
a model change. An oracle propensity option is useful for synthetic reference
calculations, but supplied true parameters cannot be counted as evidence that
an estimated-propensity workflow is calibrated.

There is also a reproducibility issue visible directly in the source. The prose
describes \texttt{integrate\_allocation} as defaulting to false, while the
inspected function signature sets it to true. The proper response is to specify
and record the argument explicitly. Section~\ref{sec:exec24-interference-reuse}
traces the downstream ratio and derivative paths; their flag propagation,
allocation-mass and missing-group restrictions still require matched native
qualification. It is not appropriate to select whichever default gives a more
favorable result after comparing outputs. \cite{v13_sw_inferference}

For large groups, multiplying many probabilities can create numerical problems.
The documentation discusses an allocation-integration option in this context.
Qualification should include large and uneven groups, boundary allocation
probabilities, integration tolerances and variance derivatives. A successful
floating-point result does not prove adequate support; an underflow error does
not prove nonidentification. These are different failure modes and need
different diagnostics.

The observed DESCRIPTION declares GPL version 2 or later and a dated version.
That informs a dependency review but does not establish the license status of
every linked component or distribution arrangement. A candidate grouped-policy
reference should remain distinct from a decision to embed or port it.

\section{A common comparison must distinguish network and group assumptions}

The two packages can be compared on a common small allocation problem only
after their target, exposure mapping and sampling assumptions are aligned.
For the two-borrower example, the exact four-assignment enumeration supplies a
transparent functional reference. A larger benchmark should vary degree,
group size, allocation probability, support, observation errors and dependence.
It should include examples where a direct-contact or group restriction is false.

Metrics should separate the target mean, direct and spillover contrasts,
coverage, finite availability, Monte Carlo integration error and total cost.
Increasing simulated allocations reduces integration error conditional on a
fitted model. It does not create independent groups or validate the graph.
Uncertainty must follow the actual sampling unit; counting nodes as independent
merely because the graph is stored in a table can produce misleading precision.

The implementation contract should retain the original graph, its direction
and weights, the identity map, intervention population, allocation rule,
background nodes and exposure mapping. A result should be rejected as
semantically incompatible if a proposed bridge silently changes any of these.
This requirement is shared by library calls, command-line jobs and agent tools.
An interface can make a scientific restriction visible; it cannot make it true.

\section{Provisional reuse choices and the evidence needed to change them}

The inspected software supports a portfolio of candidate components rather
than a single universal engine. The following choices describe potential uses and qualification conditions,
not accepted integrations or benchmarks. The earlier functional and
longitudinal chapters develop the corresponding source findings
(Chapters~\ref{v13:sw-functional} and~\ref{v13:sw-regimes}); the table retains
the conditions that must survive implementation rather than treating those
studies as still absent.

\begin{longtable}{p{0.18\textwidth}p{0.34\textwidth}p{0.39\textwidth}}
\toprule
Component & Candidate role & Conditions to preserve and qualify\\
\midrule
\endhead
DirectLiNGAM & Functional discovery adapter & Reviewed base-class behavior, prior orientation,
functional assumptions and adverse latent/Gaussian cases\\
y0 & Symbolic ID/IDC adapter & Reviewed recursive paths, conditional-path restriction,
graph validation, expression roundtrip and nonidentification references\\
delicatessen & Equation and numerical reference & Solver, Jacobian, covariance,
sampling unit and complete dependency contract\\
statsmodels & Point-treatment GMM baseline & Matched population, propensity bounds,
actual covariance behavior and dependent-data limits\\
pygformula & Longitudinal-regime adapter & History/event meaning, simulation,
bootstrap unit, diagnostics and policy scope\\
grf & R/native heterogeneous-effect component & Native algorithms, variance,
aggregation measure, ABI and license/distribution review\\
CausalGPS & Continuous-density component & Density/transformation helpers,
counterfactual doses, support and downstream estimation\\
causaldrf & Dose-response reference & Full function, interval procedures,
support, recovered holder notice and dependency terms\\
MossSpider & Restricted network-policy component & Direction, mapping, identity,
targeting, variance and observation restrictions\\
inferference & Grouped-allocation reference & Integrands, numerical defaults,
variance, group law and dependency terms\\
\bottomrule
\end{longtable}

An integration decision is stronger than identifying a promising function.
Direct adoption needs a stable, compatible dependency and scientific contract.
An adapter needs semantic parity and error preservation. A port needs the right
to reuse the implementation, an upstream behavioral oracle, independent
references, numerical/statistical equivalence criteria and a maintenance owner.
A research reference can be valuable without becoming a product dependency.
Deferral is justified when these prerequisites are missing; rejection should
state a concrete incompatibility and a revisit condition.

Commercial comparison must apply the same discipline. Vendor documentation can
establish a dated documented feature, but not measured agreement or operating
cost without an executed comparison. An inaccessible manual cannot support a
claim that the vendor lacks an operation. Our proposed graph-to-effect or
banking workflow is likewise not a delivered competitive advantage. A sound
build/buy/partner decision compares the vendor's documented and measured state
with our implemented, conditional, proposed or absent state, and makes the
missing evidence explicit.

Licensing is one part of this decision. The selected MIT and BSD texts,
repository GPL text and package license declarations provide traceable starting
points. Dependencies, bundled data, separate notices, native libraries and the
intended distribution arrangement still need review. A subprocess boundary
does not automatically settle licensing obligations, just as a Python wrapper
does not establish that a native operation has been qualified.

The software survey can therefore guide implementation immediately without
claiming premature adoption. It identifies precise operations, preserves their
mathematical meaning and specifies the evidence needed for the next decision.
Runtime, bank validation and independent scientific acceptance remain separate
from these source-grounded provisional choices.
